\documentclass[10pt,a4paper]{amsart}
\usepackage{amsmath,amssymb,amsthm,graphicx}
\usepackage{hyperref}
\numberwithin{equation}{section}

%\usepackage{subfigure}
\usepackage{upgreek}
\usepackage{mathtools}
\usepackage{tikz}
\usepackage[all]{xy}
\usetikzlibrary{matrix,calc,arrows.meta,positioning,patterns}


\newcommand{\bC}{{\mathbb C}}
\newcommand{\bP}{{\mathbb P}}
\newcommand{\bQ}{{\mathbb Q}}
\newcommand{\bR}{{\mathbb R}}
\newcommand{\bZ}{{\mathbb Z}}
\newcommand{\bN}{{\mathbb N}}
\newcommand{\bS}{{\mathbb S}}
\newcommand{\bF}{{\mathbb F}}
\newcommand{\bT}{{\mathbb T}}
\newcommand{\bK}{{\mathbb K}}
\newcommand{\bfk}{{\mathbf k}}
\newcommand{\bfq}{{\mathbf q}}


\newcommand{\cF}{\mathcal F}
\newcommand{\cG}{\mathcal G}
\newcommand{\cA}{\mathcal A}
\newcommand{\cB}{\mathcal B}
\newcommand{\cC}{\mathcal C}
\newcommand{\cD}{\mathcal D}
\newcommand{\cE}{\mathcal E}
\newcommand{\cK}{\mathcal{K}}
\newcommand{\cY}{\mathcal Y}
\newcommand{\cS}{\mathcal S}
\newcommand{\cN}{\mathcal N}
\newcommand{\cM}{\mathcal M}
\newcommand{\cO}{\mathcal O}
\newcommand{\cP}{\mathcal P}
\newcommand{\cU}{\mathcal U}
\newcommand{\cL}{\mathcal L}
\newcommand{\cR}{\mathcal R}
\newcommand{\cT}{\mathcal T}
\newcommand{\cH}{\mathcal H}
\newcommand{\oD}{\overline{D}}
\newcommand{\scrS}{\EuScript S}
\newcommand{\scrC}{\EuScript C}
\newcommand{\scrD}{\EuScript D}
\newcommand{\scrE}{\EuScript E}
\newcommand{\scrF}{\EuScript F}
\newcommand{\scrG}{\EuScript G}
\newcommand{\scrL}{\EuScript L}
\newcommand{\scrM}{\EuScript M}
\newcommand{\scrO}{\EuScript O}
\newcommand{\scrP}{\EuScript P}
\newcommand{\scrU}{\EuScript U}
\newcommand\isomto{\stackrel{\sim}{\smash{\longrightarrow}\rule{0pt}{0.4ex}}}
\newcommand{\ev}{\operatorname{ev}}
\DeclareMathOperator{\val}{val}
\DeclareMathOperator{\skel}{Skel}
\DeclareMathOperator{\id}{id}
\DeclareMathOperator{\im}{im}
\DeclareMathOperator{\hoSC}{hoSC}
\DeclareMathOperator{\whoSC}{\widehat{hoSC}}
\DeclareMathOperator{\SC}{SC}
\DeclareMathOperator{\wSC}{\widehat{SC}}
\DeclareMathOperator{\wtel}{\widehat{tel}}
\DeclareMathOperator{\Sp}{span}
\DeclareMathOperator{\tel}{tel}
\DeclareMathOperator{\inj}{inj}
\DeclareMathOperator{\conv}{conv}
\DeclareMathOperator{\Sec}{sec}
\DeclareMathOperator{\Area}{Area}
\DeclareMathOperator{\Vol}{Vol}
\DeclareMathOperator{\diam}{diam}
\DeclareMathOperator{\Tot}{Tot}
\newtheorem{Theorem}{Theorem}[section]
\newtheorem*{Theorem*}{Theorem}
\newtheorem{Corollary}[Theorem]{Corollary}
\newtheorem{Lemma}[Theorem]{Lemma}
\newtheorem{tm}[Theorem]{Theorem}
\newtheorem{lm}[Theorem]{Lemma}
\newtheorem{prp}[Theorem]{Proposition}
\newtheorem{Proposition}[Theorem]{Proposition}
\newtheorem{cy}[Theorem]{Corollary}

\theoremstyle{definition}
\newtheorem{Definition}[Theorem]{Definition}
\newtheorem{Note}[Theorem]{Note}
\newtheorem{Example}[Theorem]{Example}
\newtheorem{Remark}[Theorem]{Remark}
\newtheorem{df}[Theorem]{Definition}
\newtheorem{ques}[Theorem]{Question}
\newtheorem{ass}[Theorem]{Assumption}
\newtheorem{obs}[Theorem]{Observation}
\newtheorem{ex}[Theorem]{Example}
\newtheorem{rem}[Theorem]{Remark}
\usepackage{cleveref}
\crefname{originalreference}{original reference}{original references}
\Crefname{originalreference}{Original reference}{Original references}
\makeatletter
\newlabel{prop:starshape-hypersurface}{{1.22}{source}{}{sourcepointer-1}{}}
\newlabel{prop:starshape-hypersurface@cref}{{[originalreference][0][]1.22}{[1][0][]source}{}{}{}}
\makeatother
\title{Path connectedness and holomorphic-curve confinement for normalized complete almost complex structures}
\author{Yoel Groman}
\date{doi:10.3842/SIGMA.2026.066; CC BY 4.0}
\begin{document}
\maketitle
\noindent\textit{Allowed standard prerequisites (editorial).}\begin{itemize}
\item Symplectic cutting, holomorphic-curve compactness and standard toric/Fukaya background as used by author. Original Sikorav, STVFluxShape and FOOO/ChoOh bibliography retained; actual source-local arguments included.
\end{itemize}\medskip
\ifdefined\hypertarget\hypertarget{sourcepointer-1}{}\fi
\noindent\textit{Source reference.} Original reference 1.22: algebra-geometry/sources/sigma-2026-066/native/sigma26-066.tex, source lines 381--381. Later application or generalization mentioned in the narrative; not an input to the retained target proof.\par\medskip
\subsection{The results}
As our first proxy for completeness, we use the well-established notion of \emph{geometric boundedness} (see, e.g., \cite[Section~4.2]{Sikorav}): an~$\omega$-compatible almost complex structure $J$ on $M$ is called \emph{geometrically bounded} if the associated Riemannian metric $g_J(\cdot,\cdot)=\omega(\cdot,J\cdot)$ is complete and has uniformly bounded sectional curvature and uniformly positive injectivity radius.\footnote{To treat $\omega$-tame almost complex structures, one must also impose a~uniform bound on the norm of $\omega$ with respect to $g_J$ (i.e., $|\omega|_{g_J}\le C$).}

We begin with an existence question.
\begin{ques}\label{ques:1}
Given $\omega$, does there exist a~geometrically bounded $\omega$-compatible almost complex structure $J$ on $M$?
\end{ques}
It is customary to call $(M,\omega)$ itself geometrically bounded if the answer to Question~\ref{ques:1} is affirmative. We consider this question when we start with a~geometrically bounded $M$ and excise a~subset $N\subset M$. Question~\ref{ques:1} can be understood as the question: \emph{can the distance to $N$ of any compact subset of $M\setminus N$ be made arbitrarily large by adjusting the almost complex structure?}



\begin{tm}\label{tm:1}
Let $M$ be a~geometrically bounded symplectic manifold with finitely generated spherical homology, and suppose $\omega$ is rational.
Let $N\subset M$ be a~properly embedded symplectic hypersurface.
Then $M\setminus N$ is not geometrically bounded.
\end{tm}

\subsection{Completeness in symplectic topology}\label{subsec:IntroCompleteness}
The obstructions to geometric boundedness in Question~\ref{ques:1} come from $J$-holomorphic curves with finite energy connecting distant compact subsets. This suggests viewing geometric boundedness as a~symplectic counterpart of geodesic completeness~\cite{McDuffSymplecticStructures}, with $J$-holomorphic curves playing the role of geodesics.

However, geometric boundedness has drawbacks: it involves curvature (second derivatives), and the space of geometrically bounded structures may not be contractible or even connected. At the other extreme, one can consider the weakest completeness property.

\begin{df}[$J$-completeness]\label{df:Jcomplete}
An $\omega$-tame almost complex structure $J$ on a~symplectic manifold $(M,\omega)$ is said to be \emph{$J$-complete} if the associated Riemannian metric $g_J$ is complete and for every compact subset $K\subset M$ and every $E>0$ there exists $R>0$ such that any connected $J$-holomorphic curve whose boundary meets both $K$ and $M\setminus B_R(K)$ has energy at least $E$.
\end{df}
This is stronger than metric completeness but too weak to control the topology of the space of $J$-complete structures.

We now consider an intermediate notion between $J$-completeness and geometric boundedness which is of a~$C^0$ nature and which gives rise to a~contractible condition. It is closely related to that of intermittent boundedness which was introduced in~\cite{GromanFloerOpen}. For each $J$, we construct in Section~\ref{sec:norm-complete} an intrinsically defined continuous \emph{scale function} $\rho_J\colon M\to(0,+\infty]$ measuring the isoperimetric scale at each point. The \emph{normalized distance} $d_{\mathrm{norm},J}$ measures path lengths in units of $\rho_J$
\[
L_{\mathrm{norm},J}(\gamma):=\int_{0}^{1}\rho_J(\gamma(t))|\dot\gamma(t)|_{g_J}{\rm d}t,
\]
with $d_{\mathrm{norm},J}$ the induced path metric. The scale function depends on the choice once and for all of a~cutoff constant $c_0>\frac1{4\pi}$. The constant $c_0$ measures the allowable deviation of the metric from being flat whereas the scale function at a~point $p$ measures the size of a~neighborhood in which the metric does not deviate from flatness more than is allowed by the cutoff constant.

\begin{df}[normalized completeness]\label{df:normalized-complete-intro}
An $\omega$-compatible almost complex structure $J$ is \emph{normalized complete} if the metric space $(M,d_{\mathrm{norm},J})$ is complete.
\end{df}
\begin{rem}
In this definition and Theorem~\ref{tm:normalized-completeness}, we focus on the $\omega$-compatible case. The statements are true for the $\omega$-tame case with a~very slight adjustment. See Remark~\ref{rem:omega-tame}.
\end{rem}


\begin{tm}[basic properties of normalized completeness]\label{tm:normalized-completeness}
Let $(M,\omega)$ be a~symplectic mani\-fold.
\begin{enumerate}
\itemsep=0pt
\item[$(1)$] \emph{Connectedness.}
The space of normalized complete $\omega$-compatible almost complex structures is path connected.

\item[$(2)$] \emph{Confinement and $J$-completeness.}
If $J$ is normalized complete, then $J$ is $J$-complete {\rm(}Definition~{\rm\ref{df:Jcomplete})}.
More precisely, any connected $J$-holomorphic curve $u\colon (\Sigma,\partial\Sigma)\to(M,\omega)$ satisfies a~confinement estimate of the form
\[
u(\Sigma)\subset B^{d_{\mathrm{norm},J}}_{CE(u)+C'}(u(\partial\Sigma)),
\]
where \smash{$E(u)=\int_\Sigma u^{*}\omega$} and the constants $C,C'>0$ are universal constants.

\item[$(3)$] \emph{Equivalence to standard metric under uniform bounds.}
If $g_J$ has uniformly bounded sectional curvature and uniformly positive injectivity radius, then $d_{\mathrm{norm},J}$ and $d_{g_J}$ are bi-Lipschitz equivalent. In this case, normalized completeness is equivalent to metric completeness of $g_J$.
\end{enumerate}
\end{tm}

\section{Proofs}\label{sec:proofs}

\subsection[Energy confinement for J-holomorphic curves]{Energy confinement for $\boldsymbol{J}$-holomorphic curves}
Throughout, if $u\colon (\Sigma,\partial\Sigma)\to (M,\omega)$ is a~$J$-holomorphic curve with (possibly empty) boundary, its energy is
$
E(u):=\int_{\Sigma} u^{*}\omega \ge 0$.


The following confinement lemmas are stated for $\omega$-compatible $J$. At the end of this subsection we note that both hold for $\omega$-tame $J$ under the uniform tameness condition from the footnote to the definition of geometric boundedness on page \pageref{tm:1} (in Section 1.1).


\begin{lm}[Sikorav's confinement estimate]\label{lem:sikorav}
 Let $(M,\omega)$ be a~symplectic manifold and let $J$ be an $\omega$-compatible almost complex structure such that the associated metric $g_J(\cdot,\cdot):=\omega(\cdot,J\cdot)$ is complete and has uniformly bounded sectional curvature and positive injectivity radius $($i.e., $J$ is geometrically bounded$)$.
 Then there exist a~constant $r_0>0$ depending on the geometry and a~universal constant $c>0$ such that for every connected $J$-holomorphic curve $u$ one has
 \[
 u(\Sigma)\subset B_{r_0\,\lceil E(u)/(c r_0^2)\rceil}(u(\partial\Sigma)),
 \]
 where $B_{r}(\cdot)$ denotes the closed $g_J$-ball of radius $r$ about a~subset and $\lceil\cdot\rceil$ is the ceiling function.
 \end{lm}
 \begin{proof}
 See \cite[Proposition 4.4.1]{Sikorav}.
 \end{proof}
 This confinement estimate is concerned with the large scale. Namely, it contains no information about curves with energy much smaller than $cr_0^2$ other than that the diameter is $\leq r_0$. In some cases, e.g., in the proof of Proposition~\ref{prop:starshape-hypersurface}, we need an estimate for the small scale. This is taken care of by the following lemma.
 \begin{lm}[improved confinement]\label{lem:sikorav-improved}
 In the setting of Lemma~{\rm\ref{lem:sikorav}}, with the same $r_0$ and $c$, and letting $C := 1/c$, one has
 \[
 u(\Sigma)\subset B_{\rho}(u(\partial\Sigma)),\qquad \rho:=\max\bigl(3CE(u)/r_0,\sqrt{CE(u)}\bigr).
 \]
 \end{lm}
 \begin{proof}
 The argument relies on the same monotonicity inequality used for Sikorav's lemma: if a~$J$-holomorphic curve meets the center of a~ball of radius $r\le r_0$ and the boundary of the $J$-holomorphic curve does not meet the interior of the ball, then the energy in that ball is at least $c r^2$.

 If $E(u)/\bigl(c r_0^2\bigr)\le 1$, we apply the monotonicity inequality directly: if the image contained a~point at the center of a~ball of radius greater than \smash{$\sqrt{E(u)/c}$} whose interior did not meet $u(\partial\Sigma)$, then monotonicity would force $E(u)\ge c r^2>E(u)$ (since $E(u)\ge c r^2$ implies $r\le\sqrt{E(u)/c}$). Hence $u(\Sigma)$ is contained in a~ball of radius \smash{$\sqrt{E(u)/c} = \sqrt{C E(u)}$} about $u(\partial\Sigma)$.

 If $E(u)/\bigl(c r_0^2\bigr)>1$, take the maximum number $n$ of disjoint balls of radius $r_0$ that the image of $u$ can meet. Then $n\le E(u)/\bigl(c r_0^2\bigr)$ by monotonicity, and the image is contained in a~ball of radius $(2n+1)r_0\le r_0\bigl(1+2E(u)/\bigl(c r_0^2\bigr)\bigr)<3E(u)/(cr_0) = 3CE(u)/r_0$ around $u(\partial\Sigma)$.
 \end{proof}

For later use, we also record a~localized version, originally observed in~\cite{GromanFloerOpen}.
\begin{df}\label{df:a-bounded}
 For $a>0$, we say that the complete Riemannian metric $g_J$ is \emph{$a$-bounded at a~point} $p\in M$ if \smash{$\inj_{g_J}(p)\ge \frac{1}{a}$} and $\Sec_{g_J}\le a^{2}$ on $B_{1/a}(p)$. We say that $g_J$ is \emph{$a$-bounded on a~set} if it is $a$-bounded at every point of that set.
\end{df}
\begin{rem}
 When $g_J$ is $a$-bounded, we can take $r_0=1/(2a)$ \cite[Theorem 4.4.1 and Remark~4.4.3]{McDuffSalamonJHol}.
\end{rem}

\begin{lm}[local confinement]\label{lem:local-confinement}
Fix $a>0$. There exist $r_0=r_0(a)>0$ and a~universal constant $C>0$ with the following property. Let $K\subset M$ be compact and let $R>0$.
Assume~$J$ is $\omega$-compatible and $g_J$ is complete and is $a$-bounded on the $g_J$-ball $B_{R}(K)$.
Let $u$ be a~$J$-holomorphic curve $u$ with $u(\partial\Sigma)\subset K$ and write
\[
\rho:=\max\bigl(2CE(u)/r_0,\sqrt{CE(u)}\bigr).
\]
If $\rho\le R$, then $u(\Sigma)\subset B_R(K)$.
\end{lm}
\begin{proof}
The proof of Lemma~\ref{lem:sikorav-improved} relies on the monotonicity inequality which depends only on the local geometry of a~complete Riemannian metric. So taking $r_0$ and $C$ as in Lemma~\ref{lem:sikorav-improved}, we obtain by the same argument that
$
u(\Sigma)\subset B_{\rho}(u(\partial\Sigma))
$.
\end{proof}

\section{Normalized completeness}\label{sec:norm-complete}
\subsection{An isoperimetric scale}\label{subsec:isoperimetric-scale}
Let $(M,\omega)$ be a~symplectic manifold and let $J$ be an $\omega$-compatible almost complex structure.
Write
$
g_J(\cdot,\cdot):=\omega(\cdot,J\cdot)
$
for the associated Riemannian metric, and write $\ell_{g_J}(\gamma)$ for the $g_J$-length of a~piecewise smooth loop $\gamma\subset M$.

\begin{df}[isoperimetric ratio on a~ball]\label{df:isoperimetric-ratio}
For $p\in M$ and $r>0$, define
\[
\mathcal{I}_J(p;r)
:=
\sup_{\gamma}
\inf_{u}
\frac{\int_D u^*\omega}{\ell_{g_J}(\gamma)^2} \in[0,+\infty],
\]
where the supremum ranges over all \emph{contractible} piecewise smooth loops $\gamma\subset B^{g_J}_r(p)$ with $\ell_{g_J}(\gamma)>0$, and the infimum ranges over all \emph{smooth} discs $u\colon (D,\partial D)\to \bigl(B^{g_J}_r(p),\gamma\bigr)$ with~${\partial u=\gamma}$.
\end{df}


Fix once and for all a~cutoff constant $c_0>0$.

\begin{df}\label{df:isoperimetric-scale}
For $p\in M$, let
\[
r_{\mathrm{iso},J}(p)
:=
\sup\{r\in[0,+\infty) \mid \mathcal{I}_J(p;r)\le c_0\}
\]
and define the \emph{isoperimetric scale}
$
\rho_J(p):=\min\{1,r_{\mathrm{iso},J}(p)\}\in[0,1]$.
\end{df}
\begin{rem}\label{rem:cutoff}
The constant $c_0$ is an arbitrary fixed numerical cutoff. The effect of changing $c_0$ is discussed below.
\end{rem}

\begin{lm}[properties of the scale function]\label{lem:rho-continuous}
The function $\rho_J\colon M\to[0,1]$ is positive everywhere and $1$-Lipschitz with respect to $d_{g_J}$.
In particular, $\rho_J$ is bounded away from $0$ on compact subsets.
\end{lm}
\begin{proof}
 Positivity follows from the fact that inside a~small enough geodesic ball one has curvature bounds, hence a~quadratic isoperimetric inequality for filling contractible loops by discs; see \cite[Theorem~4.4.1]{McDuffSalamonJHol}.

Fix $p,q\in M$ and set $\delta:=d_{g_J}(p,q)$. Then for every $r>\delta$ one has the inclusion $B^{g_J}_{r-\delta}(q)\subset B^{g_J}_{r}(p)$. By Definition~\ref{df:isoperimetric-ratio}, this implies
 \[
 \mathcal{I}_J(q;r-\delta) \le \mathcal{I}_J(p;r).
 \]
 Consequently, if $\mathcal{I}_J(p;r)\le c_0$, then $\mathcal{I}_J(q;r-\delta)\le c_0$, hence
 $
 r_{\mathrm{iso},J}(q)\ge r_{\mathrm{iso},J}(p)-\delta.
 $
 By symmetry we also obtain $r_{\mathrm{iso},J}(p)\ge r_{\mathrm{iso},J}(q)-\delta$, so
 $
 |r_{\mathrm{iso},J}(p)-r_{\mathrm{iso},J}(q)|\le d_{g_J}(p,q).
 $
 This proves the claim, and hence the same for $\rho_J=\min\{1,r_{\mathrm{iso},J}\}$.
 \end{proof}

\subsection{Normalized length and normalized completeness}\label{subsec:normalized-distance}
\begin{df}[normalized length and distance]\label{df:normalized-distance}
Let $J$ be $\omega$-compatible and let $\rho_J$ be as in Definition~\ref{df:isoperimetric-scale}.
For a~piecewise smooth path $\gamma\colon [0,1]\to M$, define its \emph{normalized length} by
\[
L_{\mathrm{norm},J}(\gamma):=\int_{0}^{1}\rho_J(\gamma(t)) |\dot\gamma(t)|_{g_J}{\rm d}t.
\]
The induced path metric is denoted $d_{\mathrm{norm},J}$.
Equivalently, $d_{\mathrm{norm},J}$ is the distance associated to the conformal metric \smash{$\widetilde g_J:=\rho_J^{2}g_J$}.
\end{df}
\begin{rem}
Note that the metric $\widetilde g_J$ is generally not smooth.
\end{rem}
\begin{df}[normalized completeness]\label{df:normalized-complete}
An $\omega$-compatible almost complex structure $J$ is \emph{normalized complete} if the metric space $(M,d_{\mathrm{norm},J})$ is complete.
\end{df}

\begin{lm}[closed bounded balls are compact]\label{lem:bounded-balls-compact}
 If $J$ is normalized complete, then for any basepoint $p_0\in M$ and any $n\in\bN$, the closed $d_{\mathrm{norm},J}$-ball \smash{$K_n:=\overline B^{d_{\mathrm{norm},J}}_{n}(p_0)$} is compact in the manifold topology.
 \end{lm}
 \begin{proof}
 Since $\rho_J$ is continuous and positive, we can approximate it by a~smooth positive function~$\widetilde\rho_J$ such that the ratio $\rho_J/\widetilde\rho_J$ is uniformly bounded away from $0$ and~$+\infty$.
 The smooth Riemannian metric $\widehat g_J:=\widetilde\rho_J^{2}g_J$ is then uniformly equivalent to $\widetilde g_J=\rho_J^{2}g_J$, hence the distances~$d_{\widehat g_J}$ and~$d_{\mathrm{norm},J}$ are bi-Lipschitz equivalent.
 Since $(M,d_{\mathrm{norm},J})$ is complete, $(M,d_{\widehat g_J})$ is also complete.
 Each closed ball \smash{$K_n=\overline B^{d_{\mathrm{norm},J}}_{n}(p_0)$} is bounded with respect to $d_{\widehat g_J}$ (by uniform equivalence), and since $\widehat g_J$ is smooth, the Hopf--Rinow theorem implies that closed bounded sets are compact.
 Thus each $K_n$ is compact in the manifold topology.
 \end{proof}

\subsection{Proof of Theorem~\ref{tm:normalized-completeness}}

\subsubsection[Connectedness (Theorem~1.6 (1))]{Connectedness (Theorem~\ref{tm:normalized-completeness}\,(1))}
\begin{proof}%[Proof of Theorem~\ref{tm:normalized-completeness}(1)]
The proof of connectedness is precisely the zigzag construction in \cite[Theorem~4.7]{GromanFloerOpen}.

Denote by $\mathcal{J}^{\mathrm{comp}}_\omega$ the space of smooth $\omega$-compatible almost complex structures with the $C^\infty$ topology and by \smash{$\mathcal{J}^{\mathrm{nc,comp}}_\omega\subset \mathcal{J}^{\mathrm{comp}}_\omega$} the subspace of normalized complete structures.
Assume \smash{$\mathcal{J}^{\mathrm{nc,comp}}_\omega\neq\varnothing$}.
Without loss of generality, assume that $M$ is connected and fix a~basepoint~${p_0\in M}$.
Given $J_0,J_1\in\mathcal{J}^{\mathrm{nc,comp}}_\omega$, consider the closed normalized balls
\[
B_n^0:=\overline B^{d_{\mathrm{norm},J_0}}_{n}(p_0),\qquad B_n^1:=\overline B^{d_{\mathrm{norm},J_1}}_{n}(p_0),\qquad n\in\bN.
\]
By Lemma~\ref{lem:bounded-balls-compact}, each of these balls is compact.
By construction of the normalized distance, successive balls in each sequence are at normalized distance at least $1$
\begin{equation}\label{eq:shell-width}
d_{\mathrm{norm},J_0}\bigl( B_n^0, M\setminus B_{n+1}^0\bigr)\ge 1,\qquad d_{\mathrm{norm},J_1}\bigl( B_n^1, M\setminus B_{n+1}^1\bigr)\ge 1,\qquad n\in\bN.
\end{equation}

We choose inductively increasing subsequences $n_1^0<n_2^0<\cdots$ and $n_1^1<n_2^1<\cdots$ in $\bN$ with the following nesting property: for each $i\ge 1$,
\[
B_{n_i^0+1}^0 \subset \operatorname{Int} B_{n_i^1}^1
\qquad \text{and} \quad B_{n_i^1+1}^1 \subset \operatorname{Int}B_{n_{i+1}^0}^0.
\]
(So the ball of radius $n_i^0+1$ for $J_0$ is contained in the ball of radius $n_i^1$ for $J_1$, and the ball of radius $n_i^1+1$ for $J_1$ is contained in the ball of radius \smash{$n_{i+1}^0$} for $J_0$.)
Such subsequences exist because both sequences exhaust $M$ as $n\to\infty$.
For each $i\ge 1$, set \smash{$V_i^0:=\operatorname{Int}B_{n_{i}^0+1}^0\setminus B_{n_i^0}^0$} and \smash{$V_i^1:=\operatorname{Int}B_{n_{i}^1+1}^1\setminus B_{n_i^1}^1$}.
These are pairwise disjoint shells escaping to infinity, and each has normalized transverse width at least $1$ with respect to the corresponding $J_0$ or $J_1$ by~\eqref{eq:shell-width}.

We choose a~smooth path $J_s$, $s\in[0,1]$, in $\mathcal{J}^{\mathrm{comp}}_\omega$ from $J_0$ to $J_1$ (through $\omega$-compatible structures) which satisfies
\[
J_s\equiv J_0\quad \text{on }\bigcup_i V_i^0\ \text{for }s\in[0,2/3],
\qquad
J_s\equiv J_1\quad \text{on }\bigcup_i V_i^1\ \text{for }s\in[1/3,1].
\]


We claim that each $J_s$ is normalized complete.
Indeed, fix $s\in[0,1]$. Then at least one of the two collections of shells,
$\big\{V_i^0\big\}$ (if $s\le 2/3$) or $\big\{V_i^1\big\}$ (if $s\ge 1/3$),
consists of shells on which~$J_s$ agrees with a~normalized complete structure (either $J_0$ or $J_1$) and hence has normalized transverse width $\ge 1$.
Any $d_{\mathrm{norm},J_s}$-Cauchy sequence cannot cross infinitely many such shells (each crossing costs at least $1$), so it is eventually contained in some \smash{$B_{n_i^0}^0$} or \smash{$B_{n_i^1}^1$}, which is compact in the manifold topology.
Since $d_{\mathrm{norm},J_s}$ induces the manifold topology (it is the length metric of a~continuous conformal factor $\rho_{J_s}^{2}g_{J_s}$ with $\rho_{J_s}>0$), Cauchy sequences in a~compact set converge.
Thus $(M,d_{\mathrm{norm},J_s})$ is complete.
This shows in particular that $\mathcal{J}^{\mathrm{nc,comp}}_\omega$ is path connected.
\end{proof}

\subsubsection[Normalized confinement and J-completeness (Theorem 1.16 (2))]{Normalized confinement and $\boldsymbol{J}$-completeness (see Theorem~\ref{tm:normalized-completeness}\,(2))}
\begin{proof}

 The proof is an adjustment to Sikorav's confinement argument \cite[Proposition~4.4.1]{Sikorav}, with the bounded-geometry radius replaced pointwise by the scale function $\rho_J$ coming from the local quadratic isoperimetric inequality $\mathcal{I}_J\le c_0$.

 Fix $p\in M$ and set $r:=\rho_J(p)\in(0,1]$.
 By definition of $r_{\mathrm{iso},J}$ and $c_0$, every contractible loop~${\gamma\subset B^{g_J}_{r}(p)}$ has a~smooth filling $u$ with $\int_D u^*\omega\le c_0\,\ell_{g_J}(\gamma)^2$.
 The standard monotonicity estimate for $J$-holomorphic curves (see \cite[Proposition 4.3.1]{Sikorav}) therefore applies on the ball~$B^{g_J}_{r}(p)$ and gives: if $u$ is a~$J$-holomorphic curve whose image meets $p$ and whose boundary is disjoint from~$B^{g_J}_{r}(p)$, then
 \begin{equation}\label{eq:mono-scale}
 E\bigl(u; u^{-1}\bigl(B^{g_J}_{r}(p)\bigr)\bigr) \ge \frac{r^2}{4c_0}.
 \end{equation}

 Let $u\colon(\Sigma,\partial\Sigma)\to(M,\omega)$ be a~connected $J$-holomorphic curve, and set $A:=u(\partial\Sigma)$.
 We define the core of the image as the set of points sufficiently far from the boundary:
 \[ U := \{ x \in u(\Sigma) \mid d_{g_J}(x, A) \ge \rho_J(x) \}. \]
 For every $x \in U$, the extrinsic ball \smash{$B^{g_J}_{\rho_J(x)}(x)$} is disjoint from $A$. Consequently, applying the monotonicity estimate~\eqref{eq:mono-scale} to each such ball guarantees that the energy of the curve within it is at least $\rho_J(x)^2 / 4c_0$.

 Consider the family of extrinsic balls \smash{$\mathcal{B} = \big\{ B^{g_J}_{\rho_J(x)}(x) \big\}_{x \in U}$}. By the Vitali covering lemma, there exists a~countable subcollection
 \[
 \big\{ B_i = B^{g_J}_{a_i}(x_i) \big\}_{i \in I} , \qquad \text{with $a_i = \rho_J(x_i)$},
 \]
consisting of pairwise disjoint balls such that their $5$-times expansions \smash{$\widehat{B}_i = B^{g_J}_{5a_i}(x_i)$} completely cover $U$.
 Because the balls $B_i$ are pairwise disjoint in $M$, their preimages are disjoint in $\Sigma$. Summing the energy over this subcollection avoids any overcounting, yielding
 \[%\label{eq:energy-sum}
 \sum_{i \in I} a_i^2 \le 4c_0 \sum_{i \in I} E\bigl(u; u^{-1}(B_i)\bigr) \le 4c_0 E(u).
 \]

 We now estimate the $d_{\mathrm{norm},J}$-diameter of the expanded balls $\widehat{B}_i$. For any two points $y, z \in \widehat{B}_i$, consider a~$g_J$-geodesic $\gamma$ connecting them; its length satisfies $L \le 10a_i$. Since $\rho_J$ is $1$-Lipschitz, the value of the scale function along $\gamma$ is bounded by $\rho_J(x_i) + d_{g_J}(x_i, \gamma(t)) \le a_i + 5a_i = 6a_i$. The normalized distance is thus bounded by the integral of the scale along $\gamma$
 \[
 d_{\mathrm{norm},J}(y, z) \le \int_\gamma \rho_J(\gamma(t)) {\rm d}t \le (6a_i) \cdot L \le 60 a_i^2.
 \]
 Thus, the normalized diameter of each $\widehat{B}_i$ is at most $60a_i^2$.
 Because the normalized metric integrates the scale, the normalized diameter of \smash{$\widehat{B}_i$} is bounded by its maximum scale ($6a_i$) times its $g_J$-diameter ($10a_i$), which is at most $60a_i^2$.

 For any point $p \in u(\Sigma)$, we can now bound its normalized distance to the boundary $A$.
 If~${p \notin U}$, then $d_{g_J}(p, A) < \rho_J(p)$. A~$g_J$-shortest path from $p$ to $A$ has $g_J$-length at most $\rho_J(p)$, and the scale along this path is at most $2\rho_J(p)$. Thus, its normalized length is at most $2\rho_J(p)^2 \le 2$.
 If $p \in U$, any path in $M$ connecting $p$ to $A$ must exit $U$. We can trace a~path to the boundary of $U$ entirely within the cover $\big\{\widehat{B}_i\big\}$. The normalized distance to traverse this cover is strictly bounded by the sum of the normalized diameters of all the expanded balls. Once outside $U$, the remaining normalized distance to $A$ is at most $2$.

 We can therefore unconditionally estimate the normalized distance from any $p \in u(\Sigma)$ to the boundary $A$ by combining these bounds
 \[%\label{eq:diameter-upper-bound}
 d_{\mathrm{norm},J}(p,A) \le \sum_{i \in I} 60a_i^2 + 2 \le 240 c_0 E(u) + 2.\tag*{\qed}
 \] \renewcommand{\qed}{}
\end{proof}

In the formulation of Theorem~\ref{tm:normalized-completeness}\,(2), we only dealt with the larger energy regime. Here we record a~localized version which includes also the low-energy regime.
\begin{lm}[normalized local confinement]\label{lem:normalized-local-confinement}
Let $J$ be an $\omega$-compatible almost complex structure that is normalized complete.
There exist universal constants $C_1,C_2>0$ such that for every connected $J$-holomorphic curve $u\colon(\Sigma,\partial\Sigma)\to(M,\omega)$,
\[
u(\Sigma) \subset B^{d_{\mathrm{norm},J}}_{\rho}(u(\partial\Sigma)),\qquad \rho := \max\bigl(C_1\,E(u),C_2\,\sqrt{E(u)}\bigr).
\]
In particular, as $E(u)\to 0$ the $d_{\mathrm{norm},J}$-diameter of $u(\Sigma)$ tends to zero. Moreover, this containment is unaffected by the behavior of $J$ outside the normalized ball of radius $\rho$.
\end{lm}
\begin{proof}
We revisit the end of the proof of Theorem~\ref{tm:normalized-completeness}\,(2). Using the same notation, the normalized distance from any point $p \in u(\Sigma)$ to the boundary $A = u(\partial\Sigma)$ is at most the sum of the normalized diameters of the expanded balls $\widehat{B}_i$ covering the core $U$, plus the normalized distance from an exit point $q \notin U$ to $A$.

As observed there, the normalized diameter of the cover $\bigcup \widehat{B}_i$ is bounded by $\sum_i 60a_i^2$. By monotonicity, $\sum_i a_i^2 \le 4c_0 E(u)$, so the sum of the diameters is at most $240c_0 E(u)$.

For the final segment from $q \notin U$ to $A$, the $g_J$-distance is some length $L < \rho_J(q) \le 1$. Because the scale function is globally bounded by $1$, the normalized distance of this shortest path is at most its $g_J$-length $L$. Furthermore, the standard monotonicity estimate applied to the ball $B^{g_J}_L(q)$ gives $E(u) \ge L^2 / 4c_0$. Thus, the normalized distance from $q$ to $A$ is bounded by~$\smash{L \le \sqrt{4c_0 E(u)}}$.

Combining these yields
\[
d_{\mathrm{norm}, J}(p, A) \le 240c_0 E(u) + \sqrt{4c_0 E(u)}.
\]

If $E(u) \ge 1$, we have $\sqrt{E(u)} \le E(u)$, and we can absorb the square root term into the linear term by taking $C_1 = 240c_0 + \sqrt{4c_0}$.
If $E(u) < 1$, we have $E(u) < \sqrt{E(u)}$, and we can absorb the linear term to get an estimate $\le C_2 \sqrt{E(u)}$ by taking $C_2 = 240c_0 + \sqrt{4c_0}$.

The last sentence is immediate since the monotonicity estimate depends only on the local behavior.
\end{proof}

\subsubsection[Equivalence to standard metric under uniform bounds (Theorem 1.16 (3))]{Equivalence to standard metric under uniform bounds (Theorem~\ref{tm:normalized-completeness}\,(3))}
\begin{proof}%[Proof of Theorem~\ref{tm:normalized-completeness}(3)]
Under the assumption,
$g_J$ is $a$-bounded at every point $p\in M$. See Definition~\ref{df:a-bounded}. So the scale function satisfies $\rho_J(p)\in [2/a,1]$. See \cite[Lemma~4.10]{GromanFloerOpen}. Therefore, the scaled metric is bi-Lipschitz equivalent to the metric $g_J$.
\end{proof}

\begin{rem}[$\omega$-tame case]\label{rem:omega-tame}
If instead of $\omega$-compatible one uses an $\omega$-tame almost complex structure that is geometrically bounded and satisfies the uniform tameness condition $|\omega|_{g_J}\le C$, then one still obtains an isoperimetric scale and scale function as in Section~\ref{subsec:isoperimetric-scale}, but with a~cutoff constant $c_0'$ depending on $C$ in place of $c_0=\tfrac12$. See \cite[Section~4.4]{McDuffSalamonJHol} and \cite[Section~3, Appendix~B]{GromanWrappedSemiToricSYZ}. The statements of Theorem~\ref{tm:normalized-completeness}\,(2) and (3) then hold with the same proofs after adjusting constants.
\end{rem}

I used AI-assisted tools for editing and improving the exposition (in particular for wording, structure, and catching typographical issues).
All mathematical content, results, and any remaining errors are my own responsibility.

\begin{thebibliography}{99}
\footnotesize\itemsep=0pt

\bibitem{AurouxAnticanonicalTDuality}
Auroux D., Mirror symmetry and {$T$}-duality in the complement of an
 anticanonical divisor, \textit{J.~G\"okova Geom. Topol. GGT} \textbf{1}
 (2007), 51--91, \href{http://arxiv.org/abs/0706.3207}{arXiv:0706.3207}.

\bibitem{ChoOhToric}
Cho C.-H., Oh Y.-G., Floer cohomology and disc instantons of {L}agrangian torus
 fibers in {F}ano toric manifolds,
 \href{https://doi.org/10.4310/AJM.2006.v10.n4.a10}{\textit{Asian~J. Math.}}
 \textbf{10} (2006), 773--814,
 \href{http://arxiv.org/abs/math.SG/0308225}{arXiv:math.SG/0308225}.

\bibitem{EntovGanorMembrez}
Entov M., Ganor Y., Membrez C., Lagrangian isotopies and symplectic function
 theory, \href{https://doi.org/10.4171/CMH/451}{\textit{Comment. Math. Helv.}}
 \textbf{93} (2018), 829--882,
 \href{http://arxiv.org/abs/1610.09631}{arXiv:1610.09631}.

\bibitem{FOOO_Toric1}
Fukaya K., Oh Y.-G., Ohta H., Ono K., Lagrangian {F}loer theory on compact toric
 manifolds.~{I}, \href{https://doi.org/10.1215/00127094-2009-062}{\textit{Duke
 Math.~J.}} \textbf{151} (2010), 23--174,
 \href{http://arxiv.org/abs/0802.1703}{arXiv:0802.1703}.

\bibitem{Gallot}
Gallot S., Hulin D., Lafontaine J., Riemannian geometry, 3rd ed., \textit{Universitext},
 \href{https://doi.org/10.1007/978-3-642-18855-8}{Springer}, Berlin, 2004.

\bibitem{GromanFloerOpen}
Groman Y., Floer theory and reduced cohomology on open manifolds,
 \href{https://doi.org/10.2140/gt.2023.27.1273}{\textit{Geom. Topol.}}
 \textbf{27} (2023), 1273--1390,
 \href{http://arxiv.org/abs/1510.04265}{arXiv:1510.04265}.

\bibitem{GromanWrappedSemiToricSYZ}
Groman Y., The wrapped {F}ukaya category for semi-toric~{C}alabi--{Y}au,
 \href{https://doi.org/10.4171/jems/1466}{\textit{J.~Eur. Math. Soc. (JEMS)}}
 \textbf{26} (2024), 2373--2439,
 \href{http://arxiv.org/abs/1805.03635}{arXiv:1805.03635}.

\bibitem{GromanVarolgunesCompleteEmbeddings}
Groman Y., Varolgunes U., Locality of relative symplectic cohomology for
 complete embeddings,
 \href{https://doi.org/10.1112/s0010437x23007492}{\textit{Compos. Math.}}
 \textbf{159} (2023), 2551--2637,
 \href{http://arxiv.org/abs/2110.08891}{arXiv:2110.08891}.

\bibitem{Gunther}
G{\"u}nther P., Einige S{\"a}tze {\"u}ber das {V}olumen von geod{\"a}tischen
 {K}ugeln, \textit{Ber. Verh. S{\"a}chs. Akad. Wiss. Leipzig Math.-Nat. Kl.}
 \textbf{104} (1960), 1--38.

\bibitem{KarshonTangExcision}
Karshon Y., Tang X., Symplectic excision,
 \href{http://arxiv.org/abs/2101.03534}{arXiv:2101.03534}.

\bibitem{McDuffSymplecticStructures}
McDuff D., Symplectic structures~-- a new approach to geometry, \textit{Notices
 Amer. Math. Soc.} \textbf{45} (1998), 952--960.

\bibitem{McDuffSalamonJHol}
McDuff D., Salamon D., {$J$}-holomorphic curves and symplectic topology, 2nd ed.,
 \textit{Amer. Math. Soc. Colloq. Publ.}, Vol.~52, American
 Mathematical Society, Providence, RI, 2012.

\bibitem{STVFluxShape}
Shelukhin E., Tonkonog D., Vianna R., Geometry of symplectic flux
 and~{L}agrangian torus fibrations,
 \href{https://doi.org/10.1112/topo.70002}{\textit{J.~Topol.}} \textbf{17}
 (2024), e70002, 56~pages,
 \href{http://arxiv.org/abs/1804.02044}{arXiv:1804.02044}.

\bibitem{Sikorav}
Sikorav J.-C., Some properties of holomorphic curves in almost complex
 manifolds, in Holomorphic {C}urves in {S}ymplectic {G}eometry, \textit{Progr.
 Math.}, Vol.~117,
 \href{https://doi.org/10.1007/978-3-0348-8508-9_6}{Birkh\"auser}, Basel,
 1994, 165--189.

\bibitem{StratmannParamRays}
Stratmann B., Removing parametrized rays symplectically,
 \href{https://doi.org/10.4310/jsg.2022.v20.n2.a4}{\textit{J.~Symplectic
 Geom.}} \textbf{20} (2022), 499--508,
 \href{http://arxiv.org/abs/2009.05465}{arXiv:2009.05465}.

\end{thebibliography}
\end{document}
